\section{Laurent series}

One can also define a series for a holomorphic function around a hole, or a singularity.

\begin{defn}
Given $0 \leq r_1 < r_2 \leq \infty$ and $p \in \C$, define
\glsadd{not:annulus}%
\begin{equation*}
\ann(p;r_1,r_2)
\overset{\text{def}}{=}
\{ z \in \C : r_1 < \sabs{z - p} < r_2 \} .
\end{equation*}
When $0 < r_1 < r_2 < \infty$, we call this set an
\emph{\myindex{annulus}}\footnote{%
Q\@: What do you call a banana with a hole?
A\@: A banannulus.}.
\end{defn}

A common case is when $r_1 = 0$, that is,
the punctured disc
\begin{equation*}
\ann(p;0,r) = \Delta_r(p) \setminus \{ p \} .
\end{equation*}
When $r_2 = \infty$, on the other hand,
$\ann(p;r,\infty) = \C \setminus \overline{\Delta_{r}(p)}$ (if $r > 0$).  We will,
however, avoid the temptation of calling $\ann(p;r,\infty)$ an \myquote{annulus.}%
\footnote{\myquote{Holey plane} perhaps?  A punctured disc also ought not to be
called an \myquote{annulus,} and calling $\ann(0;0,\infty) = \C \setminus
\{ 0 \}$ an \myquote{annulus} is right out!}

\begin{thm}[Existence of Laurent series]
\index{existence of Laurent series}%
\index{Laurent series}%
\label{thm:laurent}%
Suppose that $0 \leq r_1 < r_2 \leq \infty$ and
$f \colon \ann(p;r_1,r_2) \to \C$ is holomorphic.
Then there exist unique numbers $c_n \in \C$ for $n \in \Z$ such that
\glsadd{not:laurentser}%
\begin{equation*}
f(z) = \sum_{n=-\infty}^{\infty} c_n {(z-p)}^n ,
\end{equation*}
converging uniformly absolutely on compact subsets of
$\ann(p;r_1,r_2)$.  The numbers $c_n$ are given by
\begin{equation*}
c_n = 
\frac{1}{2\pi i}
\int_{\gamma}
\frac{f(z)}{{(z-p)}^{n+1}}
\,
dz  ,
\end{equation*}
where $\gamma$ is any circle of radius $s$, $r_1 < s < r_2$, centered at
$p$, oriented counterclockwise.
\end{thm}

Recall that convergence of a double series such as
\begin{equation*}
\sum_{n=-\infty}^{\infty} a_n
\end{equation*}
means
\begin{equation*}
\sum_{n=-\infty}^{\infty} a_n
=
\lim_{N\to -\infty}
\sum_{n=N}^{-1} a_n
+
\lim_{M\to \infty}
\sum_{n=0}^{M} a_n .
\end{equation*}
That is, the limits are taken independently.
For the Laurent series,
we will generally have absolute convergence,
so the limit may be taken in any way
and in any order.
However, it is still useful to split the Laurent series into two
parts like that.
Write a Laurent series as 
\begin{multline*}
\sum_{n=-\infty}^{\infty} c_n {(z-p)}^n
=
\sum_{n=0}^{\infty} c_n {(z-p)}^n
+
\sum_{n=-\infty}^{-1} c_n {(z-p)}^n
\\
=
\sum_{n=0}^{\infty} c_n {(z-p)}^n
+
\sum_{n=1}^{\infty} c_{-n} {\left(\frac{1}{z-p}\right)}^n .
\end{multline*}
So the Laurent series behaves like two power series: one in $z-p$ and one
in $\frac{1}{z-p}$.  You can therefore apply what you know about power
series.  For example, the first one converges (uniformly absolutely on compact subsets) in $\Delta_{r_2}(p)$
for some $r_2$,
and the second one converges in $\C \setminus \overline{\Delta_{r_1}(p)}$
for some $r_1$.
If $r_1 < r_2$, then the full series converges
(uniformly absolutely on compact subsets)
in $\ann(p;r_1,r_2)$.

\begin{proof}[Proof of the theorem]
Choose two numbers $s_1$ and $s_2$ such that $r_1 < s_1 < s_2 < r_2$.
Define the cycle
\begin{equation*}
\Gamma = \partial \Delta_{s_2}(p) - \partial \Delta_{s_1}(p) .
\end{equation*}
That is, $\Gamma$ goes around the larger ($s_2$) circle counterclockwise and
around the smaller ($s_1$) circle clockwise.
See \figureref{fig:twoannuli}.

\begin{myfig}
\subimport*{figures/}{twoannuli.pdf_t}
\caption{The two annuli, the smaller annulus is shaded darker.  The two
pieces of $\Gamma$ are noted with the circular arrows.\label{fig:twoannuli}}
\end{myfig}

If $q \in \C \setminus \ann(p;r_1,r_2)$, then $n(\Gamma;q) = 0$:
If $q$ is in the \myquote{hole} of the annulus $\ann(p;r_1,r_2)$, then
$n\bigl(\partial \Delta_{s_j}(p);q\bigr) = 1$ for both $j=1,2$, and 
if $q$ is outside the annulus altogether, then
$n\bigl(\partial \Delta_{s_j}(p);q\bigr) = 0$ for both $j=1,2$ 
(see \exerciseref{exercise:windingcircle} or
\exerciseref{exercise:windingcircles}).
So $\Gamma$ is homologous to zero in the annulus $\ann(p;r_1,r_2)$.
On the other hand, if $z$ is in the (smaller) annulus $\ann(p;s_1,s_2)$,
then for similar reasons, $n(\Gamma;z) = 1$.

Via Cauchy's integral formula (\thmref{thm:CIFhomology}), for every $z \in \ann(p;s_1,s_2)$,
\begin{equation*}
f(z) = 
\frac{1}{2\pi i}
\int_{\Gamma} \frac{f(\zeta)}{\zeta-z} \, d\zeta 
=
\frac{1}{2\pi i}
\int_{\partial \Delta_{s_2}(p)} \frac{f(\zeta)}{\zeta-z} \, d\zeta 
-
\frac{1}{2\pi i}
\int_{\partial \Delta_{s_1}(p)} \frac{f(\zeta)}{\zeta-z} \, d\zeta  .
\end{equation*}

We expand the two bits separately.  First,
if $\zeta \in \partial \Delta_{s_2}$, then
$\babs{\frac{z-p}{\zeta-p}} = \frac{\sabs{z-p}}{s_2} < 1$ and so
we follow the logic of \thmref{thm:holpower}.  The reason that we can
swap the integral and the series limit is the same as in that theorem.
\begin{equation*}
\begin{split}
\frac{1}{2\pi i}
\int_{\partial \Delta_{s_2}(p)} \frac{f(\zeta)}{\zeta-z} \, d\zeta 
& =
\frac{1}{2\pi i}
\int_{\partial \Delta_{s_2}(p)} \frac{f(\zeta)}{\zeta-p}
\frac{1}{1-\frac{z-p}{\zeta-p}} \, d\zeta
\\
& =
\frac{1}{2\pi i}
\int_{\partial \Delta_{s_2}(p)} \frac{f(\zeta)}{\zeta-p}
\sum_{n=0}^\infty
{\left(\frac{z-p}{\zeta-p}\right)}^n \, d\zeta
\\
& =
\sum_{n=0}^\infty
\underbrace{
\left(
\frac{1}{2\pi i}
\int_{\partial \Delta_{s_2}(p)} \frac{f(\zeta)}{{(\zeta-p)}^{n+1}}
 \, d\zeta
\right)
}_{c_n}
{(z-p)}^n .
\end{split}
\end{equation*}

Similarly, 
if $\zeta \in \partial \Delta_{s_1}$, then
$\babs{\frac{\zeta-p}{z-p}} = \frac{s_1}{\sabs{z-p}} < 1$, and so
\begin{equation*}
\begin{split}
-\frac{1}{2\pi i}
\int_{\partial \Delta_{s_1}(p)} \frac{f(\zeta)}{\zeta-z} \, d\zeta 
& = 
\frac{1}{2\pi i}
\int_{\partial \Delta_{s_1}(p)} \frac{f(\zeta)}{z-p}
\frac{1}{1-\frac{\zeta-p}{z-p}} \, d\zeta
\\
& =
\frac{1}{2\pi i}
\int_{\partial \Delta_{s_1}(p)} \frac{f(\zeta)}{z-p}
\sum_{m=0}^\infty
{\left(\frac{\zeta-p}{z-p}\right)}^m \, d\zeta
\\
& =
\sum_{m=0}^\infty
\left(
\frac{1}{2\pi i}
\int_{\partial \Delta_{s_1}(p)} f(\zeta){(\zeta-p)}^{m}
 \, d\zeta
\right)
{(z-p)}^{-m-1}
\\
& =
\sum_{n=-\infty}^{-1}
\underbrace{
\left(
\frac{1}{2\pi i}
\int_{\partial \Delta_{s_1}(p)} \frac{f(\zeta)}{{(\zeta-p)}^{n+1}}
 \, d\zeta
\right)
}_{c_n}
{(z-p)}^{n} .
\end{split}
\end{equation*}
Adding these together we have the right thing, except that the formula for $c_n$
is not quite right.  Given any $s$ such that
$r_1 < s < r_2$,
the cycle
$\partial \Delta_{s}(p) - \partial \Delta_{s_1}(p)$ is homologous to zero in
$\ann(p;r_1,r_2)$, and 
$z \mapsto \frac{f(z)}{{(z-p)}^{n+1}}$ is holomorphic in 
$\ann(p;r_1,r_2)$.  Cauchy's theorem (\thmref{thm:CThomology}) thus says
\begin{equation*}
0 = \int_{\partial \Delta_{s}(p) - \partial \Delta_{s_1}(p)}
\frac{f(\zeta)}{{(\zeta-p)}^{n+1}} \, d\zeta
=
\int_{\partial \Delta_{s}(p)}
\frac{f(\zeta)}{{(\zeta-p)}^{n+1}} \, d\zeta
-
\int_{\partial \Delta_{s_1}(p)}
\frac{f(\zeta)}{{(\zeta-p)}^{n+1}} \, d\zeta .
\end{equation*}
Similarly for $s_2$, and hence
\begin{equation*}
c_n = \frac{1}{2\pi i}
\int_{\partial \Delta_{s}(p)} \frac{f(\zeta)}{{(\zeta-p)}^{n+1}}
 \, d\zeta .
\end{equation*}
So we get the same $c_n$ no matter which $s$ we pick.

Next, convergence.
For any $\epsilon > 0$,
the geometric series used for the first part converges uniformly
absolutely when $\babs{\frac{z-p}{\zeta-p}} = \frac{\sabs{z-p}}{s_2} \leq
1-\epsilon$.  In other words, the series converges uniformly absolutely
on compact subsets of $\Delta_{s_2}(p)$ (when $\sabs{z-p} < s_2$).
%
The geometric series used for the second part converges uniformly
absolutely when $\babs{\frac{\zeta-p}{z-p}} = \frac{s_1}{\sabs{z-p}} \leq
1-\epsilon$.  In other words, the series converges uniformly absolutely
on compact subsets of $\C \setminus \overline{\Delta_{s_1}(p)}$
(when $\sabs{z-p} > s_1$).
Hence both parts (and so the entire series) 
converge uniformly absolutely on compact
subsets of $\ann(p;s_1,s_2)$.  As $s_1$ and $s_2$ were arbitrary such that
$r_1 < s_1 < s_2 < r_2$, we get that the series
converges uniformly absolutely on compact subsets of $\ann(p;r_1,r_2)$.

Finally, uniqueness of $c_n$.  Suppose $\{ d_n \}$ is another sequence such
that
\begin{equation*}
f(z)
=
\sum_{n=-\infty}^{\infty} d_n {(z-p)}^{n} ,
\end{equation*}
converging uniformly absolutely on compact subsets of $\ann(p;r_1,r_2)$.
Then
\begin{equation*}
\begin{split}
c_m = \frac{1}{2\pi i}
\int_{\partial \Delta_{s}(p)} \frac{f(\zeta)}{{(\zeta-p)}^{m+1}}
 \, d\zeta 
& =
\frac{1}{2\pi i}
\int_{\partial \Delta_{s}(p)}
\left(\sum_{n=-\infty}^{\infty} d_n {(\zeta-p)}^{n} \right)
\frac{1}{{(\zeta-p)}^{m+1}}
 \, d\zeta 
\\
& =
\frac{1}{2\pi i}
\sum_{n=-\infty}^{\infty}
d_n
\int_{\partial \Delta_{s}(p)}
{(\zeta-p)}^{n-m-1}
 \, d\zeta 
\\
& =
d_m ,
\end{split}
\end{equation*}
as the only $n$ for which
$\int_{\partial \Delta_{s}(p)}
{(\zeta-p)}^{n-m-1}
 \, d\zeta$ is nonzero is when $n=m$, that is, when we are integrating
${(\zeta-p)}^{-1}$, in which case we get $2 \pi i$.
\end{proof}
